% ID: SEL-D0056-OA-E02
% ESTADO: validada. S057.
\textbf{a)} Completamos el cuadrado:
\[
c(x)=5+0{,}00025(x-100)^2.
\]
Las velocidades $50$ y $150$ distan lo mismo de $100$, así que ambos consumos son $\boxed{5{,}625\text{ litros}}$.

\textbf{b)} $c'(x)=0{,}0005(x-100)$. El consumo decrece en $[25,100]$ y crece en $[100,175]$.

\textbf{c)} En $100$ km/h se anula el cuadrado y se obtiene el mínimo de $5$ litros. Los puntos más alejados de $100$ son los dos extremos del dominio. El máximo, $5+0{,}00025\cdot75^2=\boxed{6{,}40625\text{ litros}}$, se alcanza tanto a $25$ como a $175$ km/h.
